\documentclass[11pt,twocolumn]{article}
% Mixed directions with LuaLaTeX and babel's bidi support: Hebrew main text,
% Arabic and Persian passages, English insertions, numbers, a table and footnotes.
\usepackage[bidi=basic,layout=tabular.lists,english]{babel}
\babelprovide[import,main]{hebrew}
\babelprovide[import]{arabic}
\babelprovide[import]{persian}
\babelfont{rm}{Libertinus Serif}
\babelfont[hebrew]{rm}{David CLM}
\babelfont[hebrew]{sf}{Frank Ruehl CLM}
\babelfont[arabic,persian]{rm}[Scale=1.1]{Amiri}
\usepackage{amsmath}
\usepackage{booktabs}
\begin{document}

\section{מבוא}
זהו מסמך ניסיון שמערבב עברית עם \foreignlanguage{english}{English text}
ומספרים כמו 2026 ו־3.14.\footnote{הערת שוליים בעברית עם \foreignlanguage{english}{a few English words}.}
המשוואה $a^2 + b^2 = c^2$ צריכה להופיע משמאל לימין.

\section{ערבית ופרסית}
\begin{otherlanguage}{arabic}
هذه فقرة عربية داخل مستند عبري، فيها رقم ٤٢ وكلمة \foreignlanguage{english}{LaTeX}.
\end{otherlanguage}

\begin{otherlanguage}{persian}
این یک پاراگراف فارسی است با عدد ۱۴۰۵.
\end{otherlanguage}

\begin{otherlanguage}{english}
\section{English section}
Left-to-right text with an embedded Hebrew word: \foreignlanguage{hebrew}{שלום},
and Arabic: \foreignlanguage{arabic}{سلام}.
\end{otherlanguage}

\section{טבלה}
\begin{tabular}{lrr}
  \toprule
  עיר & אוכלוסייה & שנה \\
  \midrule
  ירושלים & 971800 & 2022 \\
  \foreignlanguage{arabic}{القاهرة} & 10100000 & 2022 \\
  \bottomrule
\end{tabular}

\begin{itemize}
  \item פריט ראשון
  \item {\sffamily פריט שני בגופן אחר}
\end{itemize}

\end{document}
